% ============================================================
%  USER INTERFACES  (un-numbered chapter)
%  Requires in preamble:  \usepackage{graphicx}  \usepackage{float}
%  Replace each ui-*.png with your actual screenshot filename.
%  Screenshot routes:
%    ui-landing.png            -> /
%    ui-login.png              -> /auth
%    ui-patient-dashboard.png  -> /patient/dashboard
%    ui-medications.png        -> /patient/medications
%    ui-routines.png           -> /patient/routines
%    ui-chatbot.png            -> /patient/chatbot
%    ui-face-recognition.png   -> /patient/face-recognition
%    ui-memory-gallery.png     -> /patient/memory-gallery
%    ui-reports.png            -> /patient/reports (or /caregiver/reports)
%    ui-caregiver-dashboard.png-> /caregiver/dashboard
%    ui-admin-dashboard.png    -> /admin/dashboard
% ============================================================
\chapter*{User Interfaces}
\addcontentsline{toc}{chapter}{User Interfaces}
\markboth{USER INTERFACES}{USER INTERFACES}

\section*{Introduction}
MemoryCare's interface was built with one thought above everything else: it had to
be easy to use for people living with cognitive impairments. With that priority in
mind, every screen was designed the same way, using large and readable text,
high-contrast elements, clear icons, and a layout that stays consistent throughout
the app, so there is no surprise element for the users and they can navigate the app
easily on their own. The core features of MemoryCare, medication and routine
reminders, the chatbot assistant, and face recognition, are all kept within easy
reach and take as few steps as possible. The app is fully responsive across laptop
and mobile, and it supports both English and Urdu, so patients and caregivers can
each use it in whichever language they find easier. The sections below present the
main screens of the application for its three user roles: patient, caregiver, and
administrator.

\section*{Landing Page}
\begin{figure}[H]\centering\includegraphics[width=\textwidth]{ui-landing.png}
\caption{Landing Page of MemoryCare}\end{figure}
The landing page is the first interface that appears when anyone opens MemoryCare.
It leads with a clear, easy-to-read headline and a brief description of what the
system offers, reminders, the AI assistant, face recognition, and caregiver
monitoring, and it keeps \textbf{Login} and \textbf{Register} easily accessible,
since this is often the very first impression a user gets of the app. The layout
adjusts smoothly to the screen size, so whether it is opened on a phone or a laptop,
the user gets a brief introduction before deciding to sign in.

\section*{Registration and Login}
\begin{figure}[H]\centering\includegraphics[width=\textwidth]{ui-login.png}
\caption{Registration and Login Screen of MemoryCare}\end{figure}
This interface handles authentication for all roles. New accounts can be created as
either a \textbf{Patient} or a \textbf{Caregiver}, while administrator accounts are
set up separately and not through self-registration. All users sign in with their
email and password, and on a successful login each user is redirected straight to
the dashboard built for their role. The system follows role-based access control, so
no user can reach pages meant for a different type of user, which protects the
integrity of the system. Invalid credentials produce a clear error message.

\section*{Patient Dashboard}
\begin{figure}[H]\centering\includegraphics[width=\textwidth]{ui-patient-dashboard.png}
\caption{Patient Dashboard of MemoryCare}\end{figure}
The patient dashboard has been intentionally kept simple. It greets the user and
shows the medications and reminders scheduled for today, and it also offers the AI
assistant whenever help is needed. Information is shown in large, easy-to-read cards,
so the less a patient has to rely on memory, the less confusing the experience
becomes.

\section*{Medications}
\begin{figure}[H]\centering\includegraphics[width=\textwidth]{ui-medications.png}
\caption{Medications Screen of MemoryCare}\end{figure}
This interface lists the medicines the patient has been prescribed, each with its
dosage and scheduled times. When a reminder is due, the patient can mark that dose as
\textbf{taken}, which is recorded in the medication log; if a dose is missed instead,
the system flags it and sends a notification to the caregiver. Caregivers can add,
edit, or remove medications as needed, and the running log allows medication
compliance to be tracked over time.

\section*{Daily Routines}
\begin{figure}[H]\centering\includegraphics[width=\textwidth]{ui-routines.png}
\caption{Daily Routines Screen of MemoryCare}\end{figure}
This screen is dedicated to the activities that have to be completed during the day.
It tracks the patient's daily activities, for example a morning walk, meals, or
namaz, whatever they have stored as part of their routine, each with its scheduled
time. When an activity is marked complete, it is saved to the routine log, keeping
the record consistent between caregiver and patient. Routine reminders are delivered
in the same way as medication reminders, helping the patient maintain a stable daily
structure.

\section*{AI Chatbot (with Voice and Urdu)}
\begin{figure}[H]\centering\includegraphics[width=\textwidth]{ui-chatbot.png}
\caption{AI Chatbot Screen of MemoryCare}\end{figure}
The chatbot gives the patient a way to communicate with the app, either by typing or
by speaking. It runs on an intent-classification model trained specifically for this
project and replies in the patient's chosen language, English or Urdu. It is designed
to answer everyday, low-pressure questions, such as what the date is or what is coming
up next, and it keeps a record of past chats so the patient can revisit them if
needed.

\section*{Face Recognition}
\begin{figure}[H]\centering\includegraphics[width=\textwidth]{ui-face-recognition.png}
\caption{Face Recognition Screen of MemoryCare}\end{figure}
This feature helps with the difficulty patients face in recognising people. Using the
device camera, it captures a face and checks it against the faces already stored for
the patient. When a match is found, it shows the person's name, their relationship to
the patient, and the match confidence; if no match is found, it shows an
``unknown person'' message. Every recognition attempt is logged with a timestamp and
result for the caregiver's reference.

\section*{Memory Gallery}
\begin{figure}[H]\centering\includegraphics[width=\textwidth]{ui-memory-gallery.png}
\caption{Memory Gallery Screen of MemoryCare}\end{figure}
The memory gallery works like a photo-album option in the app. It is built as a grid
of images, each with a title and its relevant details, and is meant to help patients
hold on to and cherish memories of their loved ones.

\section*{Reports}
\begin{figure}[H]\centering\includegraphics[width=\textwidth]{ui-reports.png}
\caption{Reports Screen of MemoryCare}\end{figure}
The reports interface lets the user generate reports to track a patient's record in a
detailed or overall way. A report gives a summarised view of how the patient is doing,
for example medication completion, and can be exported as a \textbf{PDF} or
\textbf{Excel} file for sharing with doctors or family members.

\section*{Caregiver Dashboard}
\begin{figure}[H]\centering\includegraphics[width=\textwidth]{ui-caregiver-dashboard.png}
\caption{Caregiver Dashboard of MemoryCare}\end{figure}
This dashboard is dedicated to the caregiver. It shows the patients assigned to that
specific caregiver, and for each patient it shows their current status, medication and
routine compliance, along with links to view the patient's full details. This gives
the caregiver a clear picture of everyone assigned to them.

\section*{Administrator Dashboard}
\begin{figure}[H]\centering\includegraphics[width=\textwidth]{ui-admin-dashboard.png}
\caption{Administrator Dashboard of MemoryCare}\end{figure}
The administrator interface provides a system-wide view of the web application. From
here the administrator can manage user accounts (patients, caregivers, and
administrators), review overall system statistics, and monitor the system to make sure
the data stays accurate and secure.
